\section{Instruction Set Architecture}
\textbf{Instruction Set Architecture} (ISA) is an abstraction on the interface
between the hardware and low-level software. It abstracts the functionality
of the processor, such that
\begin{itemize}
  \item Programmers are aware of functionality available from machine code
  \item Computer Designers can design functionality independent of the hardware
\end{itemize}
\begin{figure}[h!]
  \centering
  \includegraphics[width=0.4\textwidth]{isa.jpg}
\end{figure}
\subsection{ISA Design}
We have two major design philosophies for ISA:
\vspace{2mm}

\centering
\begin{tabular}{c c}
  \toprule
  Complex Instruction Set Computer (CISC) & \textbf{Reduced Instruction Set Computer (RISC)} \\
  \midrule
  Complex operations with single instruction & Simple and small instruction set \\
  Smaller program size & Larger program size \\
  No room for hardware optimisation & Easier to build/optimise hardware\\
                                    & Burden on software to implement high-level statements\\
  \bottomrule
\end{tabular}

\vspace{2mm}
\raggedright
ISA Design comprises 5 main concepts:
\begin{enumerate}
  \item \textbf{Data Storage}: Storage Architecture
  \item \textbf{Memory Addressing Nodes}
  \item \textbf{Operations in the Instruction Set}
  \item \textbf{Instruction Formats}: Instruction Length, Fields
  \item \textbf{Encoding the Instruction Set}
\end{enumerate}

\subsubsection{Data Storage}
In data storage, we concern ourselves with storing the operands, the results,
and how the operands are specified. We have the common storage architecture
designs:
\begin{figure}[h!]
  \centering
  \includegraphics[width=0.65\textwidth]{storageArchi.jpeg}
\end{figure}
\begin{itemize}
  \item \textit{Stack Architecture}: Operands are implicitly on top of a stack
  \item \textit{Accumulator Architecture}: One operand is implicitly in a special
    register known as the accumulator (IBM 701, DEC PDP-8)
  \item \textit{General-Purpose Register Architecture}: Operands are explicit. Most common design.
      \begin{itemize}
        \item \textit{Register-Memory Architecture}: Exactly one operand is in
          memory (Motorola 68000, Intel 80386)
        \item \textit{Register-Register (Load-Store) Architecture}: Both
          operands are in registers (\textbf{MIPS}, DEC Alpha)
      \end{itemize}
  \item \textit{Memory-memory Architecture}: All operands are in memory. (DEC VAX)
\end{itemize}

\subsubsection{Memory Addressing Mode}
Given a $k$-bit address, we have an address space ($0$ to $2^{32}-1$) of size
$2^{k}$, and each memory transfer consists of one word of $n$ bits. In
register-register architecture, the memory addressing node consists of the
\textit{Memory Address Register}, \textit{Memory Data Register} which are used
by the address/data bus and the control lines to load/store data.
\paragraph{Endianness} Within a multi-byte word (say 4), the most significant
byte can be stored in the lowest address. That is, we construct the byte by
reading from memory sequentially. This is known as \textbf{big-endian}. On the
contrary, \textbf{little-endian} is the case when the least significant byte is
stored in the lowest address.
\begin{remark}
  Each byte itself is not reversed.
\end{remark}
\paragraph{Addressing Mode} For storage architectures involving explicit
operands, there are different methods of specification in the assembly
language.
\begin{example}[MIPS]
  There are 3 addressing modes in MIPS:
  \begin{itemize}
    \item \textit{Register}: Operands are specified by the register containing
    the value (\verb|add $t1, $t2, $t3|)
    \item \textit{Immediate}: Operand is specified directly in the instruction
    with its value (\verb|addi $t1, $t2, 3|)
    \item \textit{Displacement}: Operand is in memory with address calculated
    as Base $+$ Offset (\verb|lw $t1, 20($t2)|)
  \end{itemize}
\end{example}
\begin{example}[Other Addressing Modes]
  There exist other addressing modes:
  \begin{itemize}
    \item \textit{Register Indirect}: Operand is specified with a register storing
      its address
    \item \textit{Indexed/Base}: Operand is in memory with address calculated
      with Base $+$ Offset, each stored in a register
    \item \textit{Direct/Absolute}: Operand is specified directly in the instruction
      with its address
    \item \textit{Memory Indirect}: Operand is specified with a register
      containing a address to its address.
    \item \textit{Auto-increment}: Operand is written into the destination,
      then incremented
    \item \textit{Auto-decrement}: Operand is decremented, then written into
      the destination
    \item \textit{Scaled}: Operand is in memory with address calculated as Base
      $+$ Offset + Index $\times $ Scaling\_Factor
  \end{itemize}
\end{example}
\subsubsection{Operations in the Instruction Set}
Instructions are optimised based on usage frequency --- the speeds for load,
conditional branch, compare and store instructions are optimised (potentially
at the cost of other instructions).
\begin{figure}[h!]
  \centering
  \includegraphics[width=0.75\textwidth]{instructionSet.jpeg}
\end{figure}

\subsubsection{Instruction Format}
An instruction consists of the \textit{opcode} (a unique code specifying the
operation) and the \textit{operands} (any additional information needed for the
operations). The operation \textbf{designates} the type and size of the
operands.
\begin{multicols}{2}
  \textbf{Variable-length} instructions: Used in most CISC.
  \begin{itemize}
    \item Requires multi-step fetch and decode
    \item Allows for more flexible and compact instruction set
  \end{itemize}
  \textbf{Hybrid} instructions: a mix of both.\\
  \textbf{Fixed-length} instructions: Used in most RISC
  \begin{itemize}
    \item Allows for easy fetch and decode
    \item Simplifies pipelining and parallelism
    \item Instruction bits are scarce
  \end{itemize}
\end{multicols}

\subsubsection{Encoding the Instruction Set}
\begin{figure}[h!]
  \centering
  \includegraphics[width=0.6\textwidth]{encoding.jpeg}
\end{figure}
With fixed length instructions, we may adopt an \textbf{expanding opcode}
scheme. That is, the opcode has different lengths for different instructions.
\begin{example}[Maximum/Minimum Number of instructions]
  Given a $16$-bit instruction consisting type-A instructions with 6-bit
  opcodes and type-B instructions with 11-bit opcodes, the \textbf{maximum
  number of instructions} is given by maximising the prefixes for type-B,
  \[1 + \left( 2^{6}-1 \right) \times 2^{5} = 2017.\] 
  The \textbf{minimum number of instructions} is given by minimising the prefixes
  for type-B,
  \[1\times 2^{5}+2^{6}-1.\] 
  \begin{figure}[h!]
    \centering
    \includegraphics[width=0.35\textwidth]{expandingOpcode.jpeg}
  \end{figure}
\end{example}
To \textbf{design an expanding opcode} with specifications for the number of
instructions of each type, we first accommodate for the most restrictive case
(that is, the instruction type with the most operands).


\subsection{MIPS}
Microprocessor without Interlocked Pipelined Stages (MIPS) is a family of
RISC ISA\@. In summary, it adopts the stored-memory concept with the load-store
model:
\begin{itemize}
  \item \textbf{Stored-Memory Concept}: Instructions and data are stored in memory
  \item \textbf{Load-Store Architecture}: Faster processing by \textit{limiting
    memory access} through the bus, relying on registers for storage during
    execution
  \item The major \textbf{types of assembly instruction}: \textit{Memory},
    \textit{Calculation}, \textit{Control Flow}
  \item A typically assembly code structure is hence: (1) \textit{Load},
    (2) \textit{Compute}, (3) \textit{Store}
\end{itemize}
We typically have a limited number of registers in the processor. In MIPS, we
have $32$ registers, each referred to by its number or name. Registers
\textbf{do not have data types}, but each instruction assumes the data in the
register is of a certain type.
\begin{remark}
  In CS2100, we assume that we have enough registers to store all variables in
  our program.
\end{remark}
\begin{figure}[h!]
  \centering
  \includegraphics[width=0.65\textwidth]{registers.png}
\end{figure}

\subsubsection{MIPS Assembly Language}
\[\verb|op $r0, $r1, $r2|\qquad \verb|op $r0, $r1, Immd|\qquad \verb|op Immd|.\] 
Most MIPS (\textit{R-format}) operations consist three registers (1
destination, 2 sources).\ \textit{I-format} operations consist two registers (1
destination, 1 source) and an immediate value. While \textit{J-format}
operations uses only 1 immedate value.
\paragraph{Arithmetic Operations}
For arithmetic operations, The immediate ranges from $-2^{15}$ to $2^{15}-1$
(with a 16-bit 2s complement number system).
\paragraph{Logical Operations} In logical operations, the registers consist
of 32 raw bits, instead of a single 32-bit number.
\begin{figure}[h!]
  \centering
  \includegraphics[width=0.6\textwidth]{mipsBasic.jpeg}
\end{figure}

\paragraph{Memory Organisation}
In MIPS, the main memory can be viewed as a \textit{large, single-dimensional
array} of memory locations. Each location has an address, sequentially indexed
in the array (in bytes). This is known as \textit{byte addressing}.

\begin{remark}
  With the byte address, we can access a single byte (\textit{byte
  addressable}) or a single word (\textit{word addressable})
\end{remark}

MIPS consists of $32$ registers, each $32$-bit (4-byte) long. Each word contains
32 bits (4 bytes) and memory addresses are $32$-bit long.\ \textbf{Each word is
aligned in memory, that is, it begins at a byte that is a multiple of $4$.}
Hence, we have the memory instructions \verb|lw $t0, 12($t1)|,
\verb|sw $t0, 12($t1)| requiring the address be a multiple of 4. To load/store
misaligned values, we use the instructions \verb|lb| and \verb|sb| instead.

\begin{definition}[Variable Mapping]
  In C, during compilation, the compiler allocates a register for each variable used
  in the program. This process is known as \textbf{variable mapping}.
\end{definition}

\begin{example}[Large Constant Loading]
  To load a large (>16 bits) immediate, we first set the upper 16-bit:
  $\verb|lui $t0, 0xAAAA|$, then set the lower 16-bits:
  $\verb|ori $t0, $t0, 0xAAAA|$
\end{example}

\paragraph{Control Flow}
In performing general computing tasks, we are required to make decisions
(conditions) and perform iterations (loops). To alter the control flow of the program,
we have:
\begin{itemize}
  \item \textbf{Conditional} (branch) instructions
    \[\verb|bne $t0, $t1, label|.\] 
    \[\verb|beq $t0, $t1, label|.\] 
  \item \textbf{Unconditional} (jump) instructions
    \[\verb|j label|\]
\end{itemize}
\begin{remark}
  Labels are ``anchors'' in the assembly code indicating branch targets. They
  are not instructions nor variables, and \textbf{no memory is allocated}.
\end{remark}
These instructions (along with \verb|slt| and \verb|slti|) allow us to
implement any loop or conditional decision making in our programs.
\begin{example}[Array Loops]
  For instance, we may implement array loops as follows:
  \begin{figure}[h!]
    \centering
    \includegraphics[width=0.4\textwidth]{controlFlow.jpeg}
  \end{figure}\\
  Another optimisation would be to directly use a array pointer instead of an
  index in many cases.
\end{example}

\includepdf[pages=-]{MIPS_Reference.pdf}

\subsubsection{MIPS Encoding}
Instructions in MIPS are classified into \textit{Register Format},
\textit{Immediate Format}, or \textit{Jump format}, according to their
operands. \textbf{Instructions with same operand types have the same encoding.}

\paragraph{R-Format} For register formats, we have the following fields:
\begin{itemize}
  \item \verb|opcode| Partially specifies the instruction, equals 0 $\forall$ R-format instructions
  \item \verb|funct| Further specifies the instruction
  \item \verb|rs| Source Register, specifies register containing the first operand 
  \item \verb|rt| Target Register, specifies register containing the second operand 
  \item \verb|rd| Destination Register, specifies register to receive the
    result of computation
  \item \verb|shamt| Shift Amount, equals 0 $\forall$ non-shift instructions
\end{itemize}
\begin{figure}[h!]
  \centering
  \includegraphics[width=0.4\textwidth]{rFormat.jpeg}
\end{figure}

\paragraph{I-Format} For immediate formats, we have the following fields:
\begin{itemize}
  \item \verb|op| Uniquely specifies the instruction
  \item \verb|rs| Specifies the source register operand (if any)
  \item \verb|rt| Specifies the target register to receive the result
  \item \verb|immediate| Taken as a signed integer in 2s complement
    (\textit{except for logical operations})
\end{itemize}
\begin{figure}
\centering
\begin{subfigure}{.5\textwidth}
  \centering
  \includegraphics[width=.9\linewidth]{iFormat1.jpeg}
\end{subfigure}%
\begin{subfigure}{.5\textwidth}
  \centering
  \includegraphics[width=.78\linewidth]{iFormat2.jpeg}
\end{subfigure}
\end{figure}
\begin{definition}[Program Counter]
  The \textbf{program counter} (PC) is a special register keeping address of
  the next instruction to be executed in the processor.
\end{definition}
Branch instructions also use the I-format, where the immediate specifies the
\textbf{relative address of the next instruction (in number of words)} if the
condition is fulfilled.
\[PC:=\begin{cases} PC+4&\text{if branch is not taken}\\ \left( PC+4 \right)
+\left( \verb|immediate| \times 4 \right) & \text{if branch is
taken}\end{cases}.\]
\begin{remark}
  This allows us to specify a branch up to $\pm 2^{15}$ words $=\pm 2^{17}$ 
  bytes from the PC.
\end{remark}

\paragraph{J-Format} For jump formats, we have the following fields:
\begin{itemize}
  \item \verb|opcode| Uniquely specifies the instruction
  \item \verb|targetAddress| Consists the 5th to 30th MSB of the target
    instruction address.
\end{itemize}
The 4 MSBs of the target instruction address is taken from $PC+4$, and the 2
LSBs are $00$ due to word alignment of instructions. We hence form the full
32-bit target address from the specified 26 bits.

\begin{remark}
  Due to the use of the first $4$-bits of $PC+4$, jump instructions are limited to within
  the same block of 256MB.
\end{remark}

\subsubsection{MIPS Addressing Modes}
In MIPS, we have the following addressing modes:
\begin{itemize}
  \item \textit{Register Addressing}: (\verb|add, sub, and, xor, nor, slt, sll, srl|)
  \item \textit{Immediate Addressing}: (\verb|addi, andi, ori, xori, slti|)
  \item \textit{Base Addressing (Displacement Addressing)}: Operand is at the
  address obtained from the sum of a register and an immediate (\verb|lw, sw|)
  \item \textit{PC-Relative Addressing}: Address is obtained from the sum of PC
  and an immediate (\verb|beq, bne|)
  \item \textit{Pseudo-direct Addressing}: Address is given by $26$-bit instruction
    concatenated with upper $4$-bits of $PC+4$
\end{itemize}
